\section{Corona Post-Quantum Path}


\subsection{Ring-LWE Foundation}

\Corona{} is a threshold signature scheme built on the Ring Learning With Errors (Ring-LWE) problem \cite{lyubashevsky2010ideal, lux-corona-pq}. The security reduction is:

\begin{theorem}[Corona Security]
Breaking the \Corona{} $(t, n)$-threshold signature scheme with advantage $\epsilon$ implies solving Ring-LWE with dimension $d = 1024$ and modulus $q = 2^{32} - 5$ with advantage $\geq \epsilon / \text{poly}(n)$.
\end{theorem}

\subsection{Threshold Signing Protocol}

The $(t, n)$-threshold scheme requires $t = \lceil 2n/3 \rceil + 1$ shares to reconstruct a valid signature. The protocol proceeds in three rounds:

\begin{enumerate}
\item \textbf{Commitment}: Each signer $v_i$ samples noise $e_i \leftarrow \chi_\sigma$ from the discrete Gaussian distribution and broadcasts commitment $c_i = \text{Com}(e_i)$.
\item \textbf{Share}: Upon receiving $\geq t$ commitments, each signer broadcasts their partial signature $\sigma_i = s_i \cdot H_{\text{lat}}(B) + e_i$ where $s_i$ is their secret key share.
\item \textbf{Combine}: The aggregator collects $t$ valid shares and reconstructs $\sigma_{\text{thresh}} = \sum_{i \in \mathcal{T}} \lambda_i \sigma_i$ using Lagrange coefficients $\lambda_i$.
\end{enumerate}

\subsection{Signature and Key Sizes}

\begin{table}[h]
\centering
\begin{tabular}{lcc}
\toprule
\textbf{Component} & \textbf{Size (bytes)} & \textbf{Comparison to BLS} \\
\midrule
Public key & 1,952 & 40.7x larger \\
Secret key share & 3,904 & 121.9x larger \\
Signature share & 3,904 & 81.3x larger \\
Threshold signature & 3,904 & 81.3x larger \\
Commitment & 32 & 0.67x (smaller) \\
\bottomrule
\end{tabular}
\caption{Corona cryptographic sizes ($d = 1024$, $q = 2^{32} - 5$)}
\label{tab:corona-sizes}
\end{table}

The larger sizes motivate GPU acceleration: without batching, \Corona{} verification dominates consensus latency.

